\storyheader{``Hôm nay là sinh nhật của mình nhỉ?''}
{Ngày 06}
{

Tôi đi về từ cửa hàng tiện lợi, trên tay của tôi là một túi xách lớn
chứa đầy đồ hộp và nước tăng lực, đôi mắt cứ chực chờ được nhắm lại, đầu
đau như búa bổ, lưng thì đang đau nhức, tuy nhiên thì tôi lại không thể
nghỉ ngơi ngay khi về nhà vì vẫn còn deadline ở nhà. Mặc cho hôm nay
đáng ra phải là một ngày vui.

Tôi, Kamishiro Ayaka, một tiểu thuyết gia nổi tiếng, tác giả của nhiều
đầu truyện đình đám, lại chẳng hề có một cuộc sống sung sướng gì cho
cam, mỗi ngày đều là viết chương mới, mua đồ hộp, viết chương mới,
ăn,blại viết chương mới, và đi ngủ khi mà bình minh sắp ló dạng. Cứ thế
lặp đi lặp lại ngày này qua tháng nọ, nhiều đến mức mà tôi cũng sắp phát
ngán rồi.

Với số tiền đang chất đống ở nhà, tôi hoàn toàn có thể thuê trợ lý về
làm việc cùng mình, tuy nhiên thì tôi nhận ra rằng bản thân sẽ có thể
tập trung hơn khi làm việc một mình, thế nên là phương án đó chưa bao
giờ được thực hiện. Và cũng với số tiền đó, tôi hoàn toàn có thể đi du
lịch hay gì khác để mà đổi gió, nhưng tôi hiểu rõ rằng ngoài việc viết
tiểu thuyết ra, tôi cũng không có hứng thú về việc gì khác cả.

Tôi cứ chầm chậm bước đi, hiện tại đã khá khuya, tôi có thể cảm nhận gió
lạnh đang lùa vào người khiến tôi phải kéo cái khăn choàng lên một chút,
hiện tại đang là mùa đông nên thời tiết khá lạnh, từng ánh đèn đường
nhấp nháy như thể chúng đang giễu cợt cái cuộc sống tẻ nhạt này. Xung
quanh tối om như mực, có lẽ sẽ có một vài nguy hiểm trong màn đêm ấy
nhưng tôi không để tâm cho lắm, tôi cứ thế bước đi, dần đến gần cái căn
hộ rẻ tiền mà bản thân thuê.

Thượng đế ơi, nếu có thể, xin người hãy ban cho con một thứ gì đó thắp
sáng cái cuộc sống tẻ nhạt này.

Và rồi, trước mắt tôi, ngay dưới cầu thang đi lên phòng mình.

Tôi thấy một cái hộp các tông loại lớn.

Ngồi trong đó là một cô bé nhỏ nhắn, với hai cái tai mèo và một cái
đuôi. Quần áo ăn mặc rách rưới mặc cho thời tiết lạnh cóng hiện tại.

``Chị ơi.'' Cô bé ấy lên tiếng, giọng có chút run rẩy vì lạnh nhưng bằng
một cách nào ấy, tôi có thể cảm nhận thấy được sự hạnh phúc trong đó.
``Chị mang em về được không ạ? Em dễ nuôi lắm.''

*

Tôi dắt em ấy phòng mình, bên trong thì tất nhiên rồi, cực kỳ bừa bộn,
nhưng mà tất nhiên là nó không bẩn. Tôi không thể ngồi viết tiểu thuyết
trong cái nơi mà rác chất đống lên, bụi bặm khắp nơi và có một cái mùi
kinh khủng được, nên là tôi thường thuê đội dọn vệ sinh đến để dọn dẹp
cái phòng, và cũng siêng đi đổ rác nữa. Nên hiện tại vứt ngổn ngang ở
đây chỉ có mấy thứ như tạp chí, truyện tranh và ấy món đồ lặt vặt bị tôi
vứt lung tung thôi.

Tuy vậy, cô bé miêu nữ bên cạnh tôi có vẻ lại không quá để tâm đến cái
vẻ bừa bộn của căn phòng này, thay vào đó là đang nhìn cái túi xách lớn
trên tay tôi với một ánh mắt long lanh, miệng thì chảy nước dãi.

``Chỉ toàn là đồ hộp cả thôi.'' Tôi mở cái túi ra để con bé nhìn vào bên
trong. ``Nếu không chê thì em muốn ăn chứ?''

``Được ạ!?'' Cô bé nhìn lên tôi với ánh mắt lấp lánh, sau khi nhận được
cái gật đầu từ tôi thì em ấy nhảy cẫng lên. ``Vâng! Em ăn!''

Hào hứng ghê nhỉ? Tôi không biết cô bé này bị bỏ đói bao lâu nhưng nhìn
dáng vẻ đó thì hẳn bữa ăn cuối cùng của em ấy cách khá xa hiện tại, lâu
đến mức mà chỉ cần được ăn mấy cái đồ hộp rẻ tiền mà tôi mua về chống
đói qua ngày thì cũng đủ khiến em ấy rạng rỡ hẳn lên.

Tôi đem cái đống đồ hộp đến lò nướng để bỏ vào hâm, nếu tôi là người ăn
thì thường mở hộp ra hốc luôn chứ chả hâm làm gì cho mất thời gian. Với
cái số tiền dư sức mua một căn nhà đàng hoàng, hoặc dư để mua hẳn một
bữa ăn thượng hạng về nhà mà tôi vẫn chọn ở một căn hộ rẻ tiền, ăn đồ
hộp cũng rẻ tiền nốt thì căn bản tôi chả cần quan tâm hương vị thức ăn
để làm gì.

Cơ mà, hiện tại thì người ăn đống đồ hộp này sẽ là cái cô bé miêu nhân
mà tôi vừa nhặt về, với cái bộ dạng hào hứng thế kia thì lương tâm của
tôi không thể nào cho phép em ấy ăn đồ giống bản thân được.

Sau khi hâm nóng thức ăn, tôi đổ mấy thứ trong hộp ra từng dĩa, lấy máy
tính của tôi ra khỏi cái bàn nhỏ ở sát tường rồi để mấy cái dĩa đó lên.
Sau đó, tôi bế cô bé miêu nữ nhẹ một cách khó tin ấy đặt vào bên trong,
không quên đưa cho em ấy cái muỗng.

Trong lúc em ấy đang hào hứng ăn từng muỗng thì tôi cũng không quên ngắm
nghía một chút món quà sinh nhật bất ngờ của mình. Về cơ bản thì em ấy
trông cũng khá dễ thương với mái tóc màu trắng, cơ thể thì khá nhỏ nhắn,
và theo tôi thấy thì chỗ đó cũng không to lắm, nhưng cũng không gọi là
phẳng. Tuy nhỏ nhắn là vậy nhưng cô bé này có lẽ cũng đã khoảng 15 tuổi
rồi, tuy miêu nhân có vài đặc điểm của mèo nhưng về tuổi tác vẫn sẽ như
con người thôi.

Gì? Định gán mác tôi với cái đám lolicon đấy à? Xin lỗi nhưng tôi đây
mới chỉ có khoảng 20 tuổi thôi, tức là vẫn còn rất rất trẻ đấy nhé. Và
tuy không giống lắm nhưng tôi là một người học nhảy lớp, và được đánh
giá là có học lực xuất sắc đấy, cái cuộc sống hiện tại là do tôi đây đã
đạt được đam mê của mình là trở thành một tiểu thuyết gia mà thôi.

``Chủ nhân nè! Từ bây giờ em sẽ dọn phòng cho chị nhé!'' Trong lúc tôi
đang rửa mấy cái dĩa thì cô bé lớn tiếng đề nghị.

``Dọn dẹp á?''

``Vâng!'' Em ấy hào hứng đáp lời. ``Nếu căn phòng cứ bừa bộn như thế
này, chị sẽ gặp khó khăn trong việc tìm được đồ mình cần, vậy thì em sẽ
khiến nó thật ngăn nắp!''

``Thật ra cũng không khó khăn mấy...'' Tôi nhìn lướt qua căn phòng hiện
tại của mà mình đang thuê. ``Chị nhớ rõ mọi thứ đang ở đâu mà, nên về cơ
bản thế này là ngăn nắp ấy chứ?''

``Thật á? Em không tin!'' Cô bé vùng vằng. ``Vậy thì điều khiển lò sưởi
ở đâu?''

``Trong góc phòng, nằm bên dưới quyển truyện tranh.'' Tôi trả lời trong
khi vẫn tiếp tục công việc rửa chén của mình.

``Thế điều khiển máy lạnh!?'' Cô bé phía sau vẫn chưa chịu thua.

``Nằm lẫn bên trong mấy tờ báo, dưới cái tờ có bản tin bắt được tên tội
phạm cướp của ba ngày trước.''

Cứ như vậy, cô bé cứ liên tục hỏi tôi về tung tích của mấy món đồ vật
hàng ngày và tôi cứ liên tục tiết lộ vị trí của chúng hiện tại, căn bản
thì việc nhớ nơi mà mình vứt đồ lung tung cũng không khó cho lắm, nhất
là với người chả bao giờ viết mấy cái tình tiết ra để ghi nhớ mà toàn để
trong não như tôi.

Và khi tôi rửa chén xong thì cũng là lúc mà con bé bỏ cuộc. Khi tôi quay
người lại thì chợt thấy em ấy ngồi bó chân ở góc phòng, cả người toát ra
một nguồn năng lượng tiêu cực.

``Không thể nấu ăn nên em định giúp chủ nhân bằng cách dọn dẹp phòng.
Vậy mà cũng không thể làm được.'' Con bé gục đầu xuống với gương mặt
buồn bã. ``Em thật vô dụng, không thể làm điều gì, rồi chị sẽ vứt bỏ em
như người chủ cũ, rồi em sẽ lại lang thang một mình tiếp...''

Bi quan dữ. Mà tên chủ cũ vứt em ấy vì không thể làm việc nhà à?

Tôi tiến lại gần và xoa đầu cô bé đang buồn bã đó. Miệng mỉm cười trấn
an em ấy.

``Không cần lo, chị đâu nhặt em về để giúp chị làm việc nhà.'' Tôi nâng
mặt em ấy lên đối diện với bản thân. ``Chị chỉ là muốn có ai đó bầu bạn
cùng, thế thôi. Gì chứ tiền thì chị không thiếu.''

``Em sẽ cố gắng, để sau này giúp đỡ chị!'' Nghe mấy lời an ủi của tôi
xong thì từ dáng vẻ ủ rũ ban nãy, cô mèo nhỏ này ngay lập tức lấy lại sự
năng động.

Tuy tuổi đã khoảng 15 nhưng có vẻ bên trong thì em ấy chỉ như một đứa
trẻ mà thôi, dễ buồn nhưng cũng rất dễ để vui vẻ trở lại.

Chà, từ bây giờ tôi và em ấy sẽ sống cùng nhau nên có vẻ là cái nơi trú
ẩn tạm bợ này của tôi không thể đáp ứng được nhu cầu của cả hai rồi. Hồi
còn sống một mình thì về cơ bản tôi cũng chả cần nhà sang cửa rộng làm
gì nên cứ quất đại một căn phòng trọ chật hẹp và gần cửa hàng tiện lợi
cho tiện, nhưng bây giờ thì con số trong nhà đã tăng lên hai người, và
tôi cũng chẳng muốn để em ấy sống ở cái nơi như thế này được. Có vẻ như
đã đến lúc tôi động vào số tiền siêu khổng lồ trong tài khoản ngân hàng
của mình rồi.

``À phải rồi, em tên là gì nhỉ? Sống chung với nhau mà không biết tên
thì hơi bất tiện.'' Để giữ phép lịch sự tối thiểu, sau khi hỏi xong thì
tôi cũng tự giới thiệu. ``Chị tên là Kamishiro Ayaka.''

``Vâng, chủ nhân Ayaka! Em có một thỉnh cầu!'' Cô bé giơ tay lên, có vẻ
là đang cần tôi làm gì đó. ``Em muốn chị đặt một cái tên mới cho em,
được không ạ?''

Ngay trước khi câu hỏi ``Để làm gì vậy?'' buột khỏi miệng, tôi đã kịp
phanh lại. Tôi thật sự không hề biết cuộc sống trước đó của cô bé nhỏ
nhắn này như thể nào, nhưng nếu đã bị vứt bỏ trong cái bộ quần áo như
thế ở thời tiết như thế này thì hẳn tên chủ cũ cũng chẳng tốt đẹp gì cho
cam.

Với thái độ tích cực từ nãy đến giờ thì có lẽ em ấy đã vượt qua những
đau khổ trước kia rồi. Tôi đã gặp quá nhiều người để có thể đọc được cảm
xúc của những người xung quanh, và với cô bé này thì tôi không thấy có
một chút gì là ép bản thân phải vui vẻ cả.

Với một người mạnh mẽ như vậy, việc vứt bỏ cái tên trong quá khứ của
mình một cách không do dự như vậy cũng là một điều dễ hiểu thôi. Và nếu
như đã quyết tâm như thế thì tôi cũng chẳng muốn gợi lại ký ức đau khổ
ngày xưa của em ấy làm gì.

``Vậy thì...'' Tôi thoáng suy nghĩ một hồi, rồi đặt tay lên xoa đầu em
ấy. ``Tên của em sẽ là Serika nhé.''

``Vâng!'' Serika cười rạng rỡ, lộ rõ chiếc răng nanh bên trong miệng.
``Tên của em là Serika! Từ bây giờ mong chị giúp đỡ, thưa chủ nhân
Ayaka!''

``Chị cũng vậy, từ giờ mong được em giúp đỡ.''

*

Mắt tôi nheo nheo lại vì tiếp xúc quá lâu với màn hình máy tính, có vẻ
như chúng nó đang muốn được tôi thương tình mà thả lỏng để cho nhắm chặt
lại sau vài ngày thức trắng, cái lưng của tôi thì cũng trong tình thế
thảm hại không kém khi những cơn đau cứ liên tục nhói lên hành hạ cái
thể xác tàn tạ của tôi.

Nhưng tôi không quan tâm! Hiện tại, tôi cứ mặc sức để cho cơn mệt mỏi
hành hạ, mười ngón tay vẫn cứ liên tục bấm từng chữ một trên bàn phím,
mặc kệ luôn cả cơn đói cồn cào trong bụng của mình. Đơn giản bởi vì ngày
hôm nay là một ngày rất đặc biệt, khi mà ngày này hai năm trước, dưới
chân cầu thang ấy, tôi đã nhặt được cả thế giới của mình.

Khi tôi vẫn đang ngồi gõ máy tính, từ phía cửa vang lên tiếng mở, tiếp
sau đó là một người có đôi tai và cái đuôi mèo, ngoại hình khoảng 17
tuổi. Trên người đang mặc bộ đồng phục của ngôi trường cấp ba dành cho
miêu nữ khá nổi tiếng.

``Em về rồi đây... chủ nhân!'' Ngay khi nhìn thấy tôi ngồi trước máy
tính với cái quầng thâm dày cộm ở dưới mắt, Serika ngay lập tức nổi
điên. ``Chị lại làm việc quá sức nữa đấy hả!?''

``Em không hiểu, Serika à.'' Tôi đáp lại Serika đang trông như sắp cào
tôi đến nơi, trong khi mắt vẫn không hề rời khỏi màn hình máy tính.
``Hôm nay là ngày sinh nhật của em, chị không thể để cái deadline khốn
kiếp này cản bước chị dành buổi tối hôm nay với em được!''

``Oa, đáng sợ ghê.'' Serika đặt cái túi xách đựng mấy thứ em ấy vừa mua
ở siêu thị lên bàn, nhìn tôi với ánh mắt phán xét. ``Bộ tác giả nào
trong môi trường tự nhiên cũng trông như này à?''

Nói vậy nhưng em ấy cũng lấy đồ xuống nấu cơm, nói rằng không muốn làm
phiền tôi nữa. Mà thật ra là cũng có phiền lắm đâu, nói đúng hơn thì nếu
có sự hiện diện của con bé ở đây còn khiến tôi tăng thêm vài trăm công
lực ấy chứ, được nghe mấy câu cằn nhằn với giọng ngọt như mía lùi của em
ấy thì còn hiệu quả hơn nhiều so với việc hốc nước tăng lực nữa.

``Xong rồi!'' Tôi ngả người ra sau ghế, sau khi đã hoàn thành xong
chương mới. Dọn dẹp máy tính xong thì Serika cũng bưng cơm ra. Thức ăn
hôm nay nhiều hơn so với mọi ngày

``Chắc chị đói meo rồi nhỉ, chủ nhân?'' Serika ngồi xuống đối diện với
tôi. ``Mời chị.''

``Oa, món nào trông cũng ngon quá đi! Tay nghề của em đúng là đỉnh
nhất!'' Tôi lấy cái đũa của mình và chắp hai tay lại. ``Mời cả nhà ăn
cơm!''

Sau đó, vì đã ngồi viết truyện với cái bụng siêu đói và nhất quyết không
bỏ cái gì vào mồm ngoài thức ăn do Serika nấu, tất cả những món ăn trên
bàn đã nhanh chóng bay sạch trước cái mồm của tôi.

``Cảm ơn vì bữa ăn!'' Tôi ngả người ra sau ghế, vẻ mặt đầy thỏa mãn.

``Cảm ơn vì bữa ăn.'' Serika chắp hai tay lại, và sau đó bưng đống chén
dĩa đi cất.

Nếu mọi người đang muốn hỏi tại sao tôi lại làm một điều ác ôn như là để
chén dĩa cho Serika rửa trong khi hôm nay là sinh nhật con bé thì tôi
cũng sẽ nói luôn là trong hai năm sống chung, dưới sự dịu dàng và đảm
đang của em ấy thì bây giờ tôi đã hoàn toàn biến thành một đứa vô dụng
mất rồi.

Lần trước khi mà tôi định giúp thì đã làm vỡ một mớ chén dĩa, trong đó
có cả cái chén mà Serika thích nhất, thì đó là lúc tôi bàng hoàng nhận
ra bản thân đã hoàn toàn vô dụng nếu không có Serika, và con bé lần đó
cũng lo cho sốt vó rằng tôi có bị thương hay không, và cấm tôi được động
vào cái nhà bếp để tránh làm hại đến bản thân, tôi thì cũng không muốn
làm em ấy lo nên ngoan ngoãn nghe theo.

Vậy đấy, ngày xưa thì tôi có khả năng nấu nướng cực tốt, dọn nhà siêu
sạch và rửa bát siêu nhanh, nhưng cái sự toàn năng đó đã bị cô mèo nhỏ
của tôi tịch thu hết sạch, và đổi lại thì tôi sẽ là một đứa vô tích sự
ăn bám vợ. Quả là một món hời mà!

Sau khi đã rửa chén xong thì Serika cũng đến và ngồi xuống cạnh tôi, tuy
là vẫn giữ nét mặt bình thường như mọi khi nhưng với cái đuôi cứ quấn
chặt lấy chân tôi thế này thì hẳn em ấy đang muốn làm nũng rồi.

``Serika.'' Tôi cất tiếng gọi con bé, đồng thời vỗ vỗ lên đùi mình.
``Đến đây.''

``...Vâng.''

Serika ngoan ngoãn bò lên đùi tôi, rúc người và nằm co lại, cái đuôi của
con bé bỏ chân của tôi ra, thay vào đó là quấn quanh người tôi, má của
em ấy thì cứ cọ cọ vào bụng của tôi như thể đang tìm kiếm hơi ấm từ đó
vậy, và miệng của em ấy cứ phát ra mấy tiếng rên khe khẽ như một bé mèo
con đang thư giãn. Phải mấy lúc thế này thì tôi mới thấy em ấy thể hiện
giống con mèo, chứ bình thường cũng chả khác con người là mấy.

``À phải rồi.'' Tôi đang xoa đầu con bé thì chợt nhớ ra một chuyện quan
trọng. ``Hôm nay sinh nhật em, chị muốn tặng cho em cái này.''

``Gì vậy ạ?'' Serika ngồi dậy và hỏi tôi, giọng có vẻ vẫn còn chút luyến
tiếc. ``À phải rồi, em cũng có quà sinh nhật cho chị.''

Cả hai đứa bọn tôi ngay lập tức chạy ù đến nơi cất quà, rồi mang cái hộp
đã được gói lại cẩn thận chạy về lại ghế sofa.

``Chúc mừng sinh nhật, Serika!'' Tôi chìa hộp quà về phía em ấy.

Serika không nhớ được sinh nhật của bản thân, thế nên vào hai năm trước
thì tôi đã quyết định rằng ngày mà tôi nhặt em ấy về dưới cầu thang sẽ
chính là sinh nhật của em ấy. Trùng hợp thay thì ngày đó cũng là sinh
nhật của tôi luôn.

Chà, có lẽ nên gọi đó là định mệnh nhỉ?

``Em cũng vậy, chúc mừng sinh nhật, chủ nhân.''

Bọn tôi cùng mở hộp quà ra, bên trong hộp Serika tặng tôi là một chiếc
khăn choàng, nhìn tuy cũng không quá khác biệt với những chiếc khăn
choàng khác ngoài tiệm nhưng tôi biết rõ rằng đây là chiếc khăn choàng
do chính con bé may, và em ấy cũng đã học hỏi rất nhiều từ bạn bè. Đó
cũng là lý do tại sao Serika lại về muộn đến như vậy.

``Đây là... chủ nhân phiên bản nhồi bông ạ?'' Serika ngơ ngác. ``Nhưng
mà lại có vẻ như không được hoàn thiện cho lắm nhỉ?''

``À.'' Tôi gãi đầu, trong lòng đầy xấu hổ. ``Ờm thì, chị muốn làm một
con búp bê nhồi bông hình chị để tặng em, với ý nghĩa là chị sẽ luôn ở
bên em mọi lúc, nhưng tay nghề của chị không có được tốt như em, đấy là
con tốt nhất chị có thể làm rồi.''

``Không sao đâu ạ!'' Serika lắc đầu và ôm chặt lấy nó. ``Em thích món
quà này lắm!''

Tôi thở phào đầy nhẹ nhõm, khi mà món quà đầu tiên đã thành công. Tiếp
theo là đến thứ này.

``Serika.'' Tôi cất tiếng gọi em ấy, cùng lúc đó đứng lên và tiến về
phía cái tủ gần đó. ``Chị còn có thứ này muốn cho em nữa.''

Tôi mở hộc tủ ra và lấy một cái hộp nhỏ có hình vuông với tông màu tím,
bên trong chiếc hộp này là thứ mà tôi đã luôn muốn trao cho Serika từ
rất lâu.

Tôi tiến lại gần và ngồi cạnh Serika một lần nữa, lần này tôi nhìn thẳng
vào mắt của cô bé miêu nữ ở trước mặt. Người đã kéo tôi ra khỏi cái đời
sống thường nhật tẻ nhạt trước kia, khi tôi vẫn sống ở căn nhà trọ rẻ
tiền với số dư trong tài khoản cao ngất ngưởng.

``Cảm ơn em đã đồng hành cùng chị suốt quãng thời gian qua, cảm ơn vì em
đã xuất hiện dưới cầu thang nhà trọ ngày hôm ấy. Cảm ơn, vì đã đến bên
chị.'' Tôi mở chiếc hộp ra, bên trong là chiếc nhẫn cầu hôn mà tôi đã
mất rất nhiều thời gian để chọn lựa. Đây cũng là thứ tiếp theo sau ngôi
nhà này mà tôi đã tận dụng lượng ngân sách khổng lồ của mình để đạt
được.

Ngôi nhà tượng trưng cho cuộc sống của bọn tôi, là nơi mà cả hai trở về.

Và chiếc nhẫn, tượng trưng cho sự gắn kết giữa cả hai người bọn tôi từ
bây giờ.

``Serika, liệu em có bằng lòng tiếp tục ở lại đây, cùng chị viết tiếp
câu chuyện của hai ta từ giờ cho đến mãi mãi về sau không?''

``Vâng... vâng!'' Serika bật khóc, hai tay che lấy miệng mình như thể
không thể kiềm chế được cảm xúc. ``Em sẽ ở đây với chị, cùng chị đi đến
hết chặng đường còn lại, của cả hai chúng ta, thưa chủ nhân!''

``Vẫn còn gọi chị là chủ nhân sao?'' Tôi đeo nhẫn vào ngón tay của
Serika, đồng thời lau đi giọt nước mắt đang trải dài trên má con bé.
``Mối quan hệ của chúng ta bây giờ đâu phải như vậy nữa, đúng không?''

``...Chị Ayaka'' Serika nắm lấy tay tôi và áp má mình lên đó. ``Từ bây
giờ cho đến mãi về sau, xin nhờ chị giúp đỡ ạ.''

``Chị cũng vậy. Serika, sau này nhờ em giúp đỡ nhé.''

Câu chuyện của hai chúng tôi vẫn sẽ bước tiếp, có thể nó sẽ có những
chông gai trắc trở, có thể cả hai sẽ có những cơn giận hờn với nhau.

Nhưng tôi tin là sẽ ổn thôi.

Bởi vì, tôi và em ấy đã có thể tìm thấy nhau mà.

}